% ---
% id: receiver/def:setting-family
% kind: def
% title: "Setting family"
% summary: >-
%   A setting family gives, per setting, a photon-number block-diagonal POVM
%   whose joint no-click elements have the prefactor form of the multiclick
%   bound, and a flag that reads exact only where every prefactor is one by
%   construction. The high-photon weights are stated on it.
% prominence: standard
% status: accepted
% uses: [def:fock, def:clickalphabet, thm:multiclick-bound, def:bounddata, lem:collapse, lem:exactbasic]
% mentions: [cor:unconditioned-family, def:residuals]
% assumes: []
% refs: []
% provenance: newly-derived
% lean:
% notes: "Eq. eq:family-mc is the combinatorial half of thm:click-povm (iv), reproved here from eq:family-elements alone so that a setting family need not come from thm:click-povm."
% ---

The weights built on the multiclick bound
(Theorem~\ref{thm:multiclick-bound} onwards) and the residuals of
Definition~\ref{def:residuals} are stated per \emph{setting}: a receiver
without a probe has one setting, a receiver with a probe has one per declared
probe setting, and in both cases the object they are stated on is the same. This definition names that object once, so that the emitted value and
its guarantee are proved for it and each instance is established separately.
The core constructs the unconditioned receiver as an instance
(Corollary~\ref{cor:unconditioned-family}, one setting); a submodule that
constructs another, such as a receiver conditioned on its own probe with one
setting per declared probe setting, proves there that it is one.

\begin{definition}\label{def:setting-family}
Let $D\ge1$ and write $\mathcal{F}_D\coloneqq\mathcal{F}(\mathbb{C}^{D})$ for
the \tref{Fock space} of Definition~\ref{def:fock} over a $D$-dimensional
\term[The D-dimensional space whose Fock space a setting family acts on.]{channel input space},
with photon-number projectors $\Pi_n$ onto $\Symn{n}(\mathbb{C}^{D})$. A
\term[A finite set of settings, each with a photon-number block-diagonal POVM whose joint no-click elements have the prefactor form, and an exact flag.]{setting family}
of $K\ge1$ ports on $\mathcal{F}_D$ is a finite non-empty set $\Tset$ of
\emph{settings} together with, for every $t\in\Tset$:
\begin{enumerate}[label=(\alph*),leftmargin=2.6em]
\item a POVM $\{M_{S|t}\}_{S\subseteq[K]}$ on $\mathcal{F}_D$ every
  element of which commutes with the photon number, so that
  $M_{S|t}=\bigoplus_{n\ge0}\blk{M}{\eqn}_{S|t}$ with $\blk{M}{\eqn}_{S|t}$
  acting on $\Symn{n}(\mathbb{C}^{D})$. Its
  \term[For a set J of ports, the sum of the POVM elements of every outcome in which no port of J clicks.]{joint no-click elements}
  and its
  \term[The sum of the POVM elements of every outcome in which at least two ports click.]{multiclick element}
  are
  \begin{equation}\label{eq:family-elements}
    \Pvac{J}^{(t)}\coloneqq\sum_{S\subseteq[K]\setminus J}M_{S|t}
    \quad(J\subseteq[K]),\qquad
    M_{\mc|t}\coloneqq\sum_{S\in\MCK}M_{S|t},
  \end{equation}
  with $\MCK$ the \tref{multiclick set} of Definition~\ref{def:clickalphabet};
\item Hermitian operators $B^{(t)},B^{(t)}_1,\dots,B^{(t)}_K$ on
  $\mathbb{C}^{D}$ and scalars $\gamma^{(t)},\beta^{(t)}_1,\dots,
  \beta^{(t)}_K\in(0,1]$ satisfying the hypotheses of
  Theorem~\ref{thm:multiclick-bound}, namely \eqref{eq:bbkhyp} and
  $\gamma^{(t)}\le\beta^{(t)}_k$ for every $k$, such that the joint
  no-click elements have the prefactor form \eqref{eq:prefactorform},
  \begin{equation}\label{eq:family-prefactor}
    \Pvac{[K]}^{(t)}=\gamma^{(t)}\,\Fock\big(\id-B^{(t)}\big),\qquad
    \Pvac{[K]\setminus\{k\}}^{(t)}=\beta^{(t)}_k\,\Fock\big(\id-B^{(t)}_k\big);
  \end{equation}
\item a flag $\chi^{(t)}\in\{\mathrm{exact},\mathrm{general}\}$, the
  setting's
  \term[The flag of a setting, exact or general; it may read exact only where every prefactor is one by construction.]{exact flag},
  which may be $\mathrm{exact}$ only if $\gamma^{(t)}=\beta^{(t)}_k=1$ for
  every $k$ holds \emph{exactly, by the construction that produced the
  setting}, read off the declaration with no arithmetic performed, so that
  at an exact setting \eqref{eq:collapsehyp} reads $K\ge K$ and holds with
  equality by construction.
\end{enumerate}
The
\term[The admissible bound data built from a setting's prefactors and the extreme eigenvalues of its operators B.]{bound data of the setting}
are $\big(\gamma^{(t)},\{\beta^{(t)}_k\},c^{(t)},\{d^{(t)}_k\}\big)$ with
$c^{(t)}\coloneqq1-\lmax(B^{(t)})$ and $d^{(t)}_k\coloneqq1-\lmin(B^{(t)}_k)$
as in \eqref{eq:cd}; they are \tref{admissible bound data}
(Definition~\ref{def:bounddata}) by Theorem~\ref{thm:multiclick-bound}, and
$\Bnd^{(t)}$ denotes their \tref{bound function} \eqref{eq:beta}.

Two consequences are used throughout. First, the multiclick element has the
inclusion--exclusion form
\begin{equation}\label{eq:family-mc}
  M_{\mc|t}=\id+(K-1)\,\Pvac{[K]}^{(t)}-\sum_{k=1}^{K}\Pvac{[K]\setminus\{k\}}^{(t)} ,
\end{equation}
because completeness gives $M_{\mc|t}=\id-M_{\varnothing|t}-\sum_kM_{\{k\}|t}$
while \eqref{eq:family-elements} gives $M_{\varnothing|t}=\Pvac{[K]}^{(t)}$
and $M_{\{k\}|t}=\Pvac{[K]\setminus\{k\}}^{(t)}-\Pvac{[K]}^{(t)}$; so
Theorem~\ref{thm:multiclick-bound} applies to every setting, with
$\lmin(\blk{M}{\eqn}_{\mc|t})\ge\Bnd^{(t)}(n)$ for every $n\ge0$. Second,
$M_{\mc|t}$ is a sum of elements of a POVM, so hypothesis \eqref{eq:mcpos}
of Lemma~\ref{lem:exactbasic}, $M_{\mc|t}\succeq0$, holds at every setting.
\end{definition}
